\documentclass[12pt,a4paper]{article}

\input{/home/tabaka/Documents/GitHub/Defs/defs.tex}

\setcounter{zadaniaRozdzial}{6}

\begin{document}
	
	\begin{center}
		\LARGE Zadania na lekcje 6
	\end{center}
	\vspace{1.5cm}
	
	\ZbiorZadanieTwoXTwo{Rozwiąż równania sprowadzając je do postaci funkcji kwadratowej/wielomianowej:}{$2\sin^2x-\sin x = 1$}{$-\cos^2x-2\sin x = 2$}{$\tg^2x - 2\tg x = 3$}{$\frac{-2\cos^2x-\sin x + 1}{\sin^2 x}=0$}
	
	\ZbiorZadanieTwoXTwo{Rozwiąż równania poprzez grupowanie/wyciąganie przed nawias:}{$\sin^2x\cos^2x-\sin^2x-\cos x +1 =0$}{$2\sin x \cos x - \frac{\sqrt{3}}{2}\cos x + \frac{\sqrt{3}}{2} \sin x - \frac{3}{4}=0$}{$\cos^2x\sin x+\sin^2x -\sin x = 0$}{$\cos x \cdot \sin(\frac{x}{2}-\frac{\pi}{4})=\cos x$}
	
	\ZbiorZadanieTwoXTwo{Rozwiąż równania:}{$\cos(x+\frac{\pi}{3})\sin(x+\frac{\pi}{3})=\frac{1}{2}$}{$\cos(x-\frac{\pi}{6})=\sin x$}{$\sin x \cdot \tg \frac{x}{2}=1$}{$\sin^43x+\cos^43x=\frac{1}{2}$}
	
	\ZbiorZadanie{Rozwiąż równanie: $$\sin x + \sqrt{3}\cos x =1$$}
	
	\ZbiorZadanie{Rozwiąż równanie: $$|\tg x|-\tg x = \ctg x$$}
	
	\ZbiorZadanie{Rozwiąż równanie: $$\frac{1-\cos8x}{1-\tg x}=0$$}
	
	\ZbiorZadanie{Rozwiąż równanie: $$\frac{\sin x}{\cos x} + \frac{\cos x}{\sin x}=\frac{8\sin 2x}{3}$$}
	
	\newpage
	
	\begin{center}
		\LARGE Zadania domowe - lekcja 6
	\end{center}
	\vspace{1.5cm}
	
	\setcounterref{zadaniaIndex}{0}
	
	\ZbiorZadanie{Udowodnij, że poniższe równanie jest tożsamością trygonometryczną.
	$$\sin^4x+\sin^2x\cos^2x+\cos^2x = 1$$}

	\ZbiorZadanieTwoXTwo{Rozwiąż równania:}{$\sin^3x+2\sin^2x=0$}{$2\cos^2x+2\sqrt{2}\sin x=3$}{$\ctg^4x-3=2\ctg^2x$}{$2\sin^23x+\cos 3x-2=0$}
	
	\ZbiorZadanieTwoXTwo{Rozwiąż równania:}{$2\sin x \cos x - 2\sin x +\cos x = 1$}{$\ctg x = 2\sin x - \frac{1}{\sin x}$}{$\frac{\cos^2x-1}{\sin x}+\sin^3x=0$}{$-1=\frac{(\sin(x+\frac{\pi}{2})-1)(\sin(\frac{\pi}{2}-x)+1)}{\sin^2x}$}
	
	\ZbiorZadanieTwoXTwo{Rozwiąż równania:}{$\sin 2x = \sin x$}{$\sin^2x-\cos2x = 2$}{$\sin x \sin2x = \cos x \cos2x$}{$\sin^22x+4\cos2x=4$}
	
	\ZbiorZadanieTwoXTwo{Rozwiąz równania:}{$\frac{1}{2}\cos x - \frac{\sqrt{3}}{2}\sin x=1$}{$2\cos x + 3 = 4\cos\frac{x}{2}$}{$\cos x \cdot\sin 7x = \cos 3x \cdot \sin 5x$}{$\sin x + \sin 3x + \sin 5x = 0$}
	
	\ZbiorZadanie{Rozwiąż równanie: $$\cos 2x + 2 = 3|\cos x|$$}
	
	\ZbiorZadanie{Rozwiąż równanie: $$\sin x=\frac{|x|}{x}$$}
	
\end{document}